\section{Problem and objective}
\label{sec:objetivo-tecnicas}

\subsection{Problem}

The MMA has accumulated three decades of rising ozone (\cref{sec:objetivo}),
and its monitoring network produces, every hour and at 15 sites, the
concentrations of the precursors and the meteorological conditions that
govern its formation. The question this work addresses is which combination
of those conditions accompanies a day on which ozone exceeds the standard,
and how reliably such a day can be identified from them. Two obstacles make
it a multivariate problem rather than a comparison of thresholds: the
precursors are strongly correlated with one another (\ce{NO}, \ce{NO2} and
\ce{NOx} are the same emission measured three times), and the response is
rare under the current threshold, so a trivial rule is almost always right
without saying anything.

\subsection{Objective of the multivariate analysis}
\label{sec:obj:objetivo}

\textbf{Main objective.} To classify each station-day record as an
exceedance or non-exceedance of the MDA8 at the 51 ppb threshold, from the
same-day precursors and meteorological variables, and to identify which
latent conditions explain that classification.

\textit{Performance criterion:} area under the ROC curve $\geq 0.75$ on a
validation partition not seen during fitting, with sensitivity and
specificity reported separately. Accuracy is not used: under the current
65 ppb threshold, exceedances are 11.25\,\% of valid station-days, and a
constant classifier would reach 88.75\,\% accuracy without identifying the
phenomenon (\cref{sec:eda-cualitativas}). For that reason the main scenario
uses 51 ppb, the 5-year target of NOM-020-SSA1-2021 (\cref{tab:nom-limites}),
with 31.96\,\% exceedances \parencite{saito2015}.

Stations are not grouped by their mean ozone levels, which the station
explains in only 5.3\,\% of the MDA8 variance; instead they are assigned to
the seven zones of \cref{sec:contexto-estaciones}, whose spatial content is
supported by the data.

\subsection{Statistical techniques}
\label{sec:obj:tecnicas}

Interdependence techniques, which reduce the predictors to interpretable
axes, are combined with dependence techniques, which relate those axes to
exceedance:

\begin{itemize}
  \item \textbf{Principal component analysis (PCA)}: resolves the
    multicollinearity of the predictors and yields orthogonal axes with a
    physical interpretation.
  \item \textbf{Exploratory factor analysis (FA)}: separates common from
    specific variance and exposes redundancies that PCA absorbs without
    flagging. Its scores are the predictors of the classification models.
  \item \textbf{Logistic regression}: models the probability of exceedance
    and the contribution of each predictor as an odds ratio, without
    normality assumptions; it has already been applied to MDA8 exceedance
    from meteorological variables \parencite{otero2016}.
  \item \textbf{Linear discriminant analysis}: contrasts the logistic model
    under the assumptions of multivariate normality and equal covariances
    \parencite{hosmer2013}. Box-Cox left the skewness at
    $|\text{skewness}|<0.5$ (\cref{tab:sesgos-transformacion}), which
    approximates the first assumption.
  \item \textbf{Autoregressive logistic model}: incorporates the temporal
    dependence of the response, which the time-series characterization shows
    to persist after removing seasonality.
\end{itemize}

All classifiers are evaluated with the area under the ROC curve, not with
accuracy, because of the imbalance noted above \parencite{saito2015}. The
standardization of \cref{sec:transformaciones} is required: in original
units, pressure, with values near 710, would dominate distances over
\ce{CO}, whose values are below 10.

\subsection{Variables and their characteristics}
\label{sec:obj:variables}

Definitions, units and missing data are given in the data dictionary
(\cref{sec:diccionario}); \cref{tab:variables-tecnica} presents the role of
each group of variables in the objective.

\begin{table}[pos=t]
  \centering
  \caption{Role of each group of variables in the objective.}
  \label{tab:variables-tecnica}
  \footnotesize
  \begin{tabular}{@{}p{3.2cm} p{4.4cm}@{}}
    \toprule
    Group & Role \\
    \midrule
    \ce{O3}, MDA8 & Define the class; not predictors \\
    \addlinespace[2pt]
    \ce{NO}, \ce{NO2}, \ce{NO_x}, \ce{CO} & Predictors, morning window
      06--09 h \parencite{abdulwahab2005} \\
    \addlinespace[2pt]
    \ce{PM10}, \ce{PM_{2.5}}, \ce{SO2} & Predictors, daily mean \\
    \addlinespace[2pt]
    \texttt{TOUT}, \texttt{RH}, \texttt{SR}, \texttt{PRS}, \texttt{WSR},
      \texttt{viento\_u}, \texttt{viento\_v} & Predictors, daily summaries
      (maximum, total, mean, minimum) \\
    \addlinespace[2pt]
    \texttt{temporada}, \texttt{tipo\_dia}, \texttt{zona} & Categorical
      predictors (season, day type, zone), coded as $k-1$ indicators \\
    \bottomrule
  \end{tabular}
\end{table}

Precursors are summarized over the 06--09 h window because morning emissions
feed midday photochemistry; solar radiation is retained to distinguish
aerosol attenuation from a difference in light availability. Categorical
variables enter as $k-1$ indicators with an explicit reference category
(Northeast, weekday and winter); including all $k$ levels would create
perfect collinearity with the intercept. PCA and FA exclude these indicators
because the binary columns would artificially inflate the explained
variance.
